\documentclass[10pt,a4paper]{amsart}
\usepackage{amsmath,amssymb,amsthm,graphicx}
\usepackage{verbatim}
\newenvironment{reference}{\comment}{\endcomment}
%\newenvironment{reference}{}{}
\newenvironment{slogan}{\comment}{\endcomment}
\newenvironment{history}{\comment}{\endcomment}

% For commutative diagrams we use Xy-pic
\usepackage[all]{xy}

% We use 2cell for 2-commutative diagrams.
\xyoption{2cell}
% Theorem environments.
%
\theoremstyle{plain}
\newtheorem{theorem}[subsection]{Theorem}
\newtheorem{proposition}[subsection]{Proposition}
\newtheorem{lemma}[subsection]{Lemma}

\theoremstyle{definition}
\newtheorem{definition}[subsection]{Definition}
\newtheorem{example}[subsection]{Example}
\newtheorem{exercise}[subsection]{Exercise}
\newtheorem{situation}[subsection]{Situation}

\theoremstyle{remark}
\newtheorem{remark}[subsection]{Remark}
\newtheorem{remarks}[subsection]{Remarks}

\numberwithin{equation}{subsection}

% Macros
%
\def\lim{\mathop{\mathrm{lim}}\nolimits}
\def\colim{\mathop{\mathrm{colim}}\nolimits}
\def\Spec{\mathop{\mathrm{Spec}}}
\def\Hom{\mathop{\mathrm{Hom}}\nolimits}
\def\Ext{\mathop{\mathrm{Ext}}\nolimits}
\def\SheafHom{\mathop{\mathcal{H}\!\mathit{om}}\nolimits}
\def\SheafExt{\mathop{\mathcal{E}\!\mathit{xt}}\nolimits}
\def\Sch{\mathit{Sch}}
\def\Mor{\mathop{\mathrm{Mor}}\nolimits}
\def\Ob{\mathop{\mathrm{Ob}}\nolimits}
\def\Sh{\mathop{\mathit{Sh}}\nolimits}
\def\NL{\mathop{N\!L}\nolimits}
\def\CH{\mathop{\mathrm{CH}}\nolimits}
\def\proetale{{pro\text{-}\acute{e}tale}}
\def\etale{{\acute{e}tale}}
\def\QCoh{\mathit{QCoh}}
\def\Ker{\mathop{\mathrm{Ker}}}
\def\Im{\mathop{\mathrm{Im}}}
\def\Coker{\mathop{\mathrm{Coker}}}
\def\Coim{\mathop{\mathrm{Coim}}}

% Boxtimes
%
\DeclareMathSymbol{\boxtimes}{\mathbin}{AMSa}{"02}

%
% Macros for moduli stacks/spaces
%
\def\QCohstack{\mathcal{QC}\!\mathit{oh}}
\def\Cohstack{\mathcal{C}\!\mathit{oh}}
\def\Spacesstack{\mathcal{S}\!\mathit{paces}}
\def\Quotfunctor{\mathrm{Quot}}
\def\Hilbfunctor{\mathrm{Hilb}}
\def\Curvesstack{\mathcal{C}\!\mathit{urves}}
\def\Polarizedstack{\mathcal{P}\!\mathit{olarized}}
\def\Complexesstack{\mathcal{C}\!\mathit{omplexes}}
% \Pic is the operator that assigns to X its picard group, usage \Pic(X)
% \Picardstack_{X/B} denotes the Picard stack of X over B
% \Picardfunctor_{X/B} denotes the Picard functor of X over B
\def\Pic{\mathop{\mathrm{Pic}}\nolimits}
\def\Picardstack{\mathcal{P}\!\mathit{ic}}
\def\Picardfunctor{\mathrm{Pic}}
\def\Deformationcategory{\mathcal{D}\!\mathit{ef}}
\usepackage{cleveref}
\crefname{originalreference}{original reference}{original references}
\Crefname{originalreference}{Original reference}{Original references}
\makeatletter
\newlabel{lemma-independence-characters}{{0CKL}{source}{}{sourcepointer-1}{}}
\newlabel{lemma-independence-characters@cref}{{[originalreference][0][]0CKL}{[1][0][]source}{}{}{}}
\makeatother
\title{Artin-Schreier extensions: a Z/pZ Galois extension in characteristic p is generated by a root of x\textasciicircum{}p - x - c}
\author{The Stacks Project Authors}
\date{stacks-project:fields.tex:local-snapshot-20261012; GNU FDL 1.2 or later}
\begin{document}
\maketitle
\noindent\textit{Editorial note.} All delivered mathematical text is the author's own, copied from the pinned Stacks Project chapter file fields.tex; no statement or proof step was rewritten, reordered or removed.\par\medskip
\noindent\textit{Editorial note.} The counted proof is the author's complete argument: The chapter prefaces this lemma with the injective character G -> Z/pZ, sigma -> sigma(b) - b, and calls the Artin-Schreier theory a converse; it is not part of the counted unit.\par\medskip
\noindent\textit{Editorial note.} The label lemma-Artin-Schreier names the counted statement, and the proof environment immediately following it in the source is delivered verbatim as its proof.\par\medskip
\noindent\textit{Editorial note.} The counted statement is carried by the author's own lemma environment in the chapter source.\par\medskip
\noindent\textit{Allowed standard prerequisites (editorial).}\begin{itemize}
\item Same chapter: lemma-independence-characters (delivered card L48335), theorem-galois-theory (delivered card L48341, for the fixed-field step), lemma-finite-Galois (|G| = [L : K] = p for a finite Galois extension with group Z/pZ); freshman's dream in characteristic p
\end{itemize}\medskip
\ifdefined\hypertarget\hypertarget{sourcepointer-1}{}\fi
\noindent\textit{Source reference.} Original reference 0CKL: sources/stacks-fields/pinned-original-source.tex, source lines 1676--1676. Prerequisite pointer invoked at 3560; it names a standard fact of the same chapter that the counted argument uses, and every step of the counted proof built on it is retained verbatim in the delivered body.\par\medskip
\begin{lemma}[Artin-Schreier extensions]
\label{lemma-Artin-Schreier}
Let $L/K$ be a Galois extension of fields of characteristic $p > 0$
with Galois group $\mathbf{Z}/p\mathbf{Z}$. Then $L = K[z]$ with
$z^p - z \in K$.
\end{lemma}

\begin{proof}
Let $\sigma$ be a generator of $\text{Gal}(L/K)$.
Consider $\sigma : L \to L$ as a $K$-linear operator.
Observe that $\sigma^p - 1 = 0$ as a linear operator.
Applying linear independence of characters
(Lemma \ref{lemma-independence-characters}),
there cannot be a polynomial of degree $< p$
annihilating $\sigma$. We conclude that the minimal
polynomial of $\sigma$ is $x^p - 1 = (x - 1)^p$.
This implies that there exists $w \in L$ such that
$(\sigma - 1)^{p - 1}(w) = y$ is nonzero. Then
$\sigma(y) = y$, i.e., $y \in K$. Thus
$z = y^{-1}(\sigma - 1)^{p - 2}(w)$ satisfies
$\sigma(z) = z + 1$. Since $z \not \in K$ we have $L = K[z]$.
Moreover since $\sigma(z^p - z) = (z + 1)^p - (z + 1) = z^p - z$
we see that $z^p - z \in K$ and the proof is complete.
\end{proof}

\end{document}
